% Generated from project outputs. Do not edit by hand.
\begin{table}[!htbp]
\centering
\begin{threeparttable}
\caption{Labor income by subsample: 2017 moments}
\label{tab:annual-2017-y1-moments}
\small
\setlength{\tabcolsep}{4pt}
\renewcommand{\arraystretch}{1.15}
\begin{tabularx}{\linewidth}{@{}Xrrrrr@{}}
\toprule
Sample & N & Mean & SD & Min & Max \\
\midrule
General & 94 & 2,483.8 & 1,596.9 & 154.2 & 6,385.6 \\
Men & 94 & 2,765.4 & 1,839.1 & 207.9 & 9,273.8 \\
Women & 94 & 2,041.8 & 1,472.6 & 112.4 & 6,251.9 \\
Urban & 94 & 2,703.0 & 1,520.3 & 323.9 & 6,385.6 \\
Rural & 92 & 2,058.4 & 1,983.8 & 32.1 & 14,062.8 \\
Indigenous & 93 & 1,890.9 & 1,432.8 & 51.9 & 6,400.0 \\
Non-Indigenous & 94 & 2,717.4 & 1,551.1 & 247.1 & 6,697.7 \\
Bottom 99 percent & 94 & 2,315.0 & 1,363.1 & 154.2 & 5,931.8 \\
Bottom 90 percent & 94 & 1,894.7 & 940.5 & 122.3 & 3,717.5 \\
Bottom 50 percent & 92 & 965.4 & 398.5 & 60.5 & 2,423.9 \\
\bottomrule
\end{tabularx}
\begin{tablenotes}[flushleft]
\footnotesize
\item Monthly quetzales. Unweighted distribution of survey-weighted cohort means before transition matching. N counts cohort cells, not individuals. SD is dispersion, not a standard error. Subsamples overlap. Cumulative income definitions and treatment of missing components follow the merging scripts.
\end{tablenotes}
\end{threeparttable}
\end{table}
